\documentclass[11pt,a4paper]{article}
\usepackage{get-uebung}

\blatt{8}{Quellen, Quellenumwandlung und Arbeitspunkt}{Kap.~2.8--2.9}

% Raster fuer U-I-Diagramme (I horizontal, U vertikal wie im Skript)
\pgfplotsset{kennlinie/.style={
  width=0.62\textwidth, height=0.5\textwidth,
  axis lines=left, grid=both, minor tick num=1,
  grid style={gray!25}, major grid style={gray!45},
  xlabel={$I$ in A}, ylabel={$U$ in V},
  /pgf/number format/use comma,
  every axis x label/.style={at={(current axis.right of origin)},anchor=west},
  every axis y label/.style={at={(current axis.above origin)},anchor=south},
  every axis plot/.append style={thick}}}

\begin{document}

\blattkopf{
  \item das Verhalten idealer und realer Spannungs- und Stromquellen beschreiben und ihre Kennlinien zeichnen,
  \item Leerlaufspannung, Kurzschlussstrom und Innenwiderstand einer realen Quelle aus Messwerten bestimmen,
  \item reale Spannungs- und Stromquellen ineinander umwandeln und damit Netzwerke schrittweise vereinfachen,
  \item den Arbeitspunkt einer Quelle mit linearer und nichtlinearer Last grafisch ermitteln.
}

% -------------------------------------------------------------------
\begin{aufgabe}{Eine Autobatterie vermessen}{15}
Eine Autobatterie soll durch eine reale Spannungsquelle (Leerlaufspannung $U_0$, Innenwiderstand $R_i$) beschrieben werden. Dazu wird ein einstellbarer Lastwiderstand $R_L$ angeschlossen und es werden zwei Messpunkte aufgenommen:
\begin{center}
\begin{minipage}{0.48\textwidth}
\centering
\begin{circuitikz}
  \draw (0,0) to[V, v<=$U_0$] (0,3) to[R, l=$R_i$] (2.6,3)
        to[short, -o, i=$I$] (3.6,3) -- (5,3) to[vR, l=$R_L$] (5,0) -- (3.6,0)
        to[short, o-] (0,0);
  \draw (3.6,3) to[open, v^=$U$] (3.6,0);
  \draw[dashed, gray] (-1.4,-0.35) rectangle (2.85,3.9);
  \node[gray, font=\small] at (0.7,-0.65) {Batterie};
\end{circuitikz}
\end{minipage}\hfill
\begin{minipage}{0.45\textwidth}
\centering
\begin{tabular}{ccc}
\toprule
Messung & $I$ & $U$\\\midrule
1 & $\SI{2}{A}$ & $\SI{11.4}{V}$\\
2 & $\SI{6}{A}$ & $\SI{10.2}{V}$\\\bottomrule
\end{tabular}
\end{minipage}
\end{center}
\begin{teilaufgaben}
  \item Stellen Sie für beide Messpunkte die Gleichung $U = U_0 - R_i\cdot I$ auf und bestimmen Sie daraus $R_i$ und $U_0$.
  \item Wie groß ist der Kurzschlussstrom $I_k$? Warum sollte man ihn bei einer Autobatterie besser nicht messen?
  \item Geben Sie das Ersatzschaltbild als reale \emph{Stromquelle} an (mit Zahlenwerten).
  \item Welche Klemmenspannung und welcher Strom stellen sich bei $R_L = \SI{1.2}{\ohm}$ ein?
  \item Begründen Sie, warum das Modell der \emph{idealen} Spannungsquelle für einen Kurzschluss keinen Sinn ergibt.
\end{teilaufgaben}
\end{aufgabe}

\begin{loesung}
\schritt{a) Zwei Gleichungen, zwei Unbekannte}
\leitfrage{Beide Messpunkte liegen auf derselben Geraden. Was erhalten Sie, wenn Sie die beiden Gleichungen voneinander abziehen?}
\[
  \SI{11.4}{V} = U_0 - R_i\cdot\SI{2}{A}, \qquad \SI{10.2}{V} = U_0 - R_i\cdot\SI{6}{A}
\]
Subtraktion der zweiten von der ersten Gleichung eliminiert $U_0$:
\[
  \SI{1.2}{V} = R_i\cdot\SI{4}{A} \quad\Rightarrow\quad R_i = \frac{\Delta U}{\Delta I} = \frac{\SI{1.2}{V}}{\SI{4}{A}} = \ergebnis{\SI{0.3}{\ohm}}
\]
Einsetzen in die erste Gleichung:
\[
  U_0 = \SI{11.4}{V} + \SI{0.3}{\ohm}\cdot\SI{2}{A} = \ergebnis{\SI{12.0}{V}}
\]
Der Innenwiderstand ist also (bis auf das Vorzeichen) die Steigung der Kennlinie.
\schritt{b) Kurzschlussstrom}
Im Kurzschluss ist $U = 0$:
\[
  I_k = \frac{U_0}{R_i} = \frac{\SI{12}{V}}{\SI{0.3}{\ohm}} = \ergebnis{\SI{40}{A}}
\]
Bei einer echten Autobatterie ist $R_i$ noch deutlich kleiner (einige m$\Omega$), der Kurzschlussstrom läge im Bereich mehrerer hundert Ampere. Leitungen glühen durch, die Batterie kann beschädigt werden oder Gas freisetzen. Deshalb bestimmt man $I_k$ rechnerisch aus zwei ungefährlichen Messpunkten.
\schritt{c) Ersatzstromquelle}
Bei der Quellenumwandlung bleibt der Innenwiderstand gleich, er liegt jetzt \emph{parallel} zur idealen Stromquelle:
\[
  I_k = \SI{40}{A}, \qquad R_i = \SI{0.3}{\ohm} \ \text{parallel}
\]
\schritt{d) Belastung mit $R_L = \SI{1.2}{\ohm}$}
$R_i$ und $R_L$ liegen in Reihe:
\[
  I = \frac{U_0}{R_i + R_L} = \frac{\SI{12}{V}}{\SI{1.5}{\ohm}} = \ergebnis{\SI{8}{A}}, \qquad
  U = R_L\cdot I = \SI{1.2}{\ohm}\cdot\SI{8}{A} = \ergebnis{\SI{9.6}{V}}
\]
Kontrolle mit der Kennlinie: $U = \SI{12}{V} - \SI{0.3}{\ohm}\cdot\SI{8}{A} = \SI{9.6}{V}$ \checkmark
\schritt{e) Grenze des idealen Modells}
Eine ideale Spannungsquelle erzwingt an ihren Klemmen immer die Spannung $U_0$. Ein idealer Kurzschluss erzwingt die Spannung $0$. Beide Forderungen widersprechen sich; formal ergäbe sich $I = U_0/0 \to \infty$. Erst der Innenwiderstand macht das Modell realistisch.
\typfehler{$R_i = U/I$ aus \emph{einem} Messpunkt berechnen -- das liefert den Lastwiderstand ($\SI{11.4}{V}/\SI{2}{A} = \SI{5.7}{\ohm} = R_L$), nicht den Innenwiderstand! $R_i$ steckt in der \emph{Änderung} $\Delta U/\Delta I$.}
\end{loesung}

% -------------------------------------------------------------------
\begin{aufgabe}{Netzwerk durch Quellenumwandlung vereinfachen}{15}
Das folgende Netzwerk soll bezüglich der Klemmen a--b durch schrittweise Quellenumwandlung in eine einzige reale Spannungsquelle umgeformt werden.
\begin{center}
\begin{circuitikz}
  \draw (0,0) to[V, v<=$U$] (0,3) to[R, l=$R_1$] (3,3)
        to[R, l=$R_2$, *-*] (3,0);
  \draw (3,3) -- (5,3);
  \draw (5,0) to[I, i=$I_q$, *-*] (5,3);
  \draw (5,3) to[R, l=$R_3$] (8,3) to[short, -o] (9,3) node[right]{a};
  \draw (0,0) -- (8,0) to[short, -o] (9,0) node[right]{b};
  \draw (9,3) to[open, v^=$U_\text{ab}$] (9,0);
\end{circuitikz}
\end{center}
\gegeben{$U = \SI{24}{V}$, $I_q = \SI{2}{A}$, $R_1 = \SI{6}{\ohm}$, $R_2 = \SI{12}{\ohm}$, $R_3 = \SI{2}{\ohm}$}
\begin{teilaufgaben}
  \item Wandeln Sie die reale Spannungsquelle ($U$, $R_1$) in eine reale Stromquelle um. Zeichnen Sie die neue Schaltung.
  \item Fassen Sie die parallelen Stromquellen und die parallelen Widerstände jeweils zusammen.
  \item Wandeln Sie das Ergebnis wieder in eine reale Spannungsquelle um und berücksichtigen Sie $R_3$. Geben Sie die Leerlaufspannung $U_0$ und den Innenwiderstand $R_i$ des Netzwerks an.
  \item An die Klemmen wird $R_L = \SI{6}{\ohm}$ angeschlossen. Berechnen Sie $U_\text{ab}$ und den Strom durch $R_L$.
\end{teilaufgaben}
\begin{tipp}Bei der Umwandlung bleibt der Innenwiderstand gleich, die Quelle ändert ihren Wert nach dem Ohmschen Gesetz. Achten Sie auf die Pfeilrichtung: Der Strompfeil der Ersatzstromquelle zeigt zum Pluspol der Spannungsquelle.\end{tipp}
\end{aufgabe}

\begin{loesung}
\schritt{a) Spannungsquelle $\to$ Stromquelle}
\leitfrage{Welcher Strom würde fließen, wenn man die Klemmen der realen Spannungsquelle ($U$, $R_1$) kurzschließt?}
\[
  I_{k1} = \frac{U}{R_1} = \frac{\SI{24}{V}}{\SI{6}{\ohm}} = \SI{4}{A}, \qquad R_1 = \SI{6}{\ohm} \text{ jetzt parallel}
\]
Der Pluspol von $U$ liegt oben, also speist $I_{k1}$ ebenso wie $I_q$ nach oben in den oberen Knoten ein:
\begin{center}
\begin{circuitikz}[scale=0.8, transform shape]
  \draw (0,0) to[I, i=$I_{k1}$] (0,3) -- (6,3);
  \draw (1.8,3) to[R, l=$R_1$, *-*] (1.8,0);
  \draw (3.6,3) to[R, l=$R_2$, *-*] (3.6,0);
  \draw (5.4,0) to[I, i=$I_q$] (5.4,3);
  \draw (5.4,3) to[R, l=$R_3$, *-] (8,3) to[short, -o] (8.8,3) node[right]{a};
  \draw (0,0) -- (8,0) to[short, -o] (8.8,0) node[right]{b};
\end{circuitikz}
\end{center}
\schritt{b) Zusammenfassen}
Parallele Stromquellen addieren sich (gleiche Richtung!), parallele Widerstände werden parallel geschaltet:
\[
  I_\text{ges} = I_{k1} + I_q = \SI{4}{A} + \SI{2}{A} = \SI{6}{A}, \qquad
  R_1\parallel R_2 = \frac{\SI{6}{\ohm}\cdot\SI{12}{\ohm}}{\SI{18}{\ohm}} = \SI{4}{\ohm}
\]
\schritt{c) Stromquelle $\to$ Spannungsquelle, $R_3$ in Reihe}
\[
  U' = I_\text{ges}\cdot (R_1\parallel R_2) = \SI{6}{A}\cdot\SI{4}{\ohm} = \SI{24}{V}
\]
Jetzt liegt $R_3$ in Reihe zum Innenwiderstand $\SI{4}{\ohm}$:
\begin{center}
\begin{circuitikz}[scale=0.8, transform shape]
  \draw (0,0) to[V, v<=$U_0$] (0,3) to[R, l=$R_i$] (3,3) to[short, -o] (4,3) node[right]{a};
  \draw (0,0) -- (3,0) to[short, -o] (4,0) node[right]{b};
\end{circuitikz}
\end{center}
\[
  U_0 = \ergebnis{\SI{24}{V}}, \qquad R_i = \SI{4}{\ohm} + R_3 = \ergebnis{\SI{6}{\ohm}}
\]
Durch $R_3$ fließt im Leerlauf kein Strom, deshalb ändert er die Leerlaufspannung nicht.
\schritt{d) Belastung}
Mit $R_L = \SI{6}{\ohm} = R_i$ teilt sich die Spannung genau zur Hälfte auf:
\[
  U_\text{ab} = U_0\cdot\frac{R_L}{R_i+R_L} = \SI{24}{V}\cdot\frac{6}{12} = \ergebnis{\SI{12}{V}}, \qquad
  I_L = \frac{U_\text{ab}}{R_L} = \ergebnis{\SI{2}{A}}
\]
\typfehler{Den Widerstand $R_3$ bei der Umwandlung „mitnehmen“, obwohl er nicht direkt parallel bzw.\ in Reihe zur jeweiligen Quelle liegt. Umwandeln darf man nur eine Quelle mit \emph{ihrem} Innenwiderstand. Außerdem: Die Pfeilrichtung der neuen Stromquelle falsch wählen -- dann subtrahieren sich die Ströme in b) fälschlich zu $\SI{2}{A}$.}
\end{loesung}

% -------------------------------------------------------------------
\begin{aufgabe}{Arbeitspunkt mit ohmscher Last grafisch bestimmen}{12}
Eine reale Spannungsquelle mit $U_0 = \SI{12}{V}$ und $R_i = \SI{4}{\ohm}$ wird mit einem Lastwiderstand $R_L$ belastet.
\begin{teilaufgaben}
  \item Berechnen Sie den Kurzschlussstrom und zeichnen Sie die Quellenkennlinie $U(I)$ in das Diagramm.
  \item Zeichnen Sie die Widerstandsgeraden für $R_L = \SI{8}{\ohm}$ und $R_L = \SI{2}{\ohm}$ ein und lesen Sie die Arbeitspunkte ab.
  \item Überprüfen Sie einen der Arbeitspunkte rechnerisch.
  \item Welcher Lastwiderstand führt zum Arbeitspunkt $I = \SI{1.5}{A}$? Welche Klemmenspannung stellt sich ein?
\end{teilaufgaben}
\begin{center}
\begin{tikzpicture}
\begin{axis}[kennlinie, xmin=0, xmax=3.5, ymin=0, ymax=13, xtick={0,0.5,...,3.5}, ytick={0,2,...,12}]
  \ifloesung
    \addplot[hsublau, domain=0:3] {12-4*x} node[pos=0.1, above right, font=\small]{Quelle};
    \addplot[loesungfarbe, domain=0:1.625] {8*x} node[pos=0.9, left, font=\small]{$\SI{8}{\ohm}$};
    \addplot[hinweisfarbe, domain=0:3.5] {2*x} node[pos=0.95, above left, font=\small]{$\SI{2}{\ohm}$};
    \addplot[only marks, mark=*, mark size=2.5pt] coordinates {(1,8) (2,4)};
    \node[right, font=\small] at (axis cs:1,8) {AP$_1$};
    \node[above right, font=\small] at (axis cs:2,4) {AP$_2$};
  \fi
\end{axis}
\end{tikzpicture}
\end{center}
\end{aufgabe}

\begin{loesung}
\schritt{a) Quellenkennlinie}
Die Kennlinie $U = U_0 - R_i\cdot I$ ist eine Gerade durch die beiden Achsenschnittpunkte
\[
  (I = 0;\ U = U_0 = \SI{12}{V}) \quad\text{und}\quad (I = I_k = U_0/R_i = \SI{3}{A};\ U = 0).
\]
\schritt{b) Widerstandsgeraden und Arbeitspunkte}
\leitfrage{Welche Kurve beschreibt $U = R_L\cdot I$? Durch welchen Punkt geht sie immer?}
Jede Widerstandsgerade geht durch den Ursprung; einen zweiten Punkt erhält man z.\,B.\ bei $I = \SI{1}{A}$: $U = \SI{8}{V}$ bzw.\ $U = \SI{2}{V}$. Der Schnittpunkt mit der Quellenkennlinie ist der Arbeitspunkt (siehe Diagramm oben):
\[
  R_L = \SI{8}{\ohm}:\ \ergebnis{I = \SI{1}{A},\ U = \SI{8}{V}}
  \qquad\qquad
  R_L = \SI{2}{\ohm}:\ \ergebnis{I = \SI{2}{A},\ U = \SI{4}{V}}
\]
\schritt{c) Rechnerische Kontrolle}
Für $R_L = \SI{8}{\ohm}$: $I = U_0/(R_i+R_L) = \SI{12}{V}/\SI{12}{\ohm} = \SI{1}{A}$, $U = \SI{8}{\ohm}\cdot\SI{1}{A} = \SI{8}{V}$ \checkmark
\schritt{d) Gesuchter Lastwiderstand}
Auf der Quellenkennlinie gilt bei $I = \SI{1.5}{A}$: $U = \SI{12}{V} - \SI{4}{\ohm}\cdot\SI{1.5}{A} = \SI{6}{V}$. Damit
\[
  R_L = \frac{U}{I} = \frac{\SI{6}{V}}{\SI{1.5}{A}} = \ergebnis{\SI{4}{\ohm}}
\]
Hier ist $R_L = R_i$, der Arbeitspunkt liegt genau in der Mitte der Kennlinie ($U_0/2$, $I_k/2$). In Kapitel 2.12 werden Sie sehen, dass die Quelle in diesem Punkt die maximale Leistung abgibt.
\typfehler{Die Widerstandsgerade mit der Steigung $1/R_L$ zeichnen. Weil hier $U$ über $I$ aufgetragen ist, ist die Steigung $R_L$ selbst. Prüfen Sie immer einen Punkt mit $U = R_L\cdot I$ nach.}
\end{loesung}

% -------------------------------------------------------------------
\begin{aufgabe}{Arbeitspunkt mit nichtlinearem Bauteil}{15}
Dieselbe Quelle wie in Aufgabe 8.3 ($U_0 = \SI{12}{V}$, $R_i = \SI{4}{\ohm}$) versorgt jetzt ein nichtlineares Bauteil N. Seine Kennlinie wurde gemessen und ist unten eingezeichnet:
\begin{center}
\begin{tabular}{l*{6}{c}}
\toprule
$U$ in V & 0 & 2 & 4 & 6 & 8 & 10\\
$I$ in A & 0 & 0{,}17 & 0{,}67 & 1{,}50 & 2{,}67 & 4{,}17\\\bottomrule
\end{tabular}
\end{center}
\begin{teilaufgaben}
  \item Warum kann man den Strom hier nicht einfach mit dem Spannungsteiler oder $I = U_0/(R_i + R)$ berechnen?
  \item Bestimmen Sie den Arbeitspunkt grafisch.
  \item Welchen Widerstand $R = U/I$ hat das Bauteil im Arbeitspunkt? Ist es ein Widerspruch, dass Sie diesen Wert jetzt doch angeben können?
  \item Die Leerlaufspannung der Quelle wird auf $U_0' = \SI{16}{V}$ erhöht ($R_i$ unverändert). Zeichnen Sie die neue Quellenkennlinie ein und lesen Sie den neuen Arbeitspunkt ab. Um wie viel Prozent steigt der Strom?
\end{teilaufgaben}
\begin{center}
\begin{tikzpicture}
\begin{axis}[kennlinie, xmin=0, xmax=4.5, ymin=0, ymax=17, xtick={0,0.5,...,4.5}, ytick={0,2,...,16}]
  \addplot[black, domain=0:4.5, samples=60] {sqrt(24*x)} node[pos=0.97, below right, font=\small]{N};
  \ifloesung
    \addplot[hsublau, domain=0:3] {12-4*x} node[pos=0.12, right=2pt, font=\small]{$U_0 = \SI{12}{V}$};
    \addplot[loesungfarbe, domain=0:4] {16-4*x} node[pos=0.12, right=2pt, font=\small]{$U_0' = \SI{16}{V}$};
    \addplot[only marks, mark=*, mark size=2.5pt] coordinates {(1.5,6) (2.188,7.247)};
    \node[below left, font=\small] at (axis cs:1.5,6) {AP};
    \node[below left, font=\small] at (axis cs:2.188,7.247) {AP$'$};
  \fi
\end{axis}
\end{tikzpicture}
\end{center}
\end{aufgabe}

\begin{loesung}
\schritt{a) Warum keine Formel?}
Spannungsteiler und $I = U_0/(R_i+R)$ setzen einen \emph{konstanten} Widerstand voraus. Bei N hängt das Verhältnis $U/I$ vom Strom ab (die Kennlinie ist keine Gerade durch den Ursprung). Man kennt $R$ also erst, wenn man den Arbeitspunkt schon kennt.
\schritt{b) Grafische Lösung}
\leitfrage{Welche beiden Bedingungen muss der Arbeitspunkt gleichzeitig erfüllen?}
Der Strom durch die Quelle ist derselbe wie durch N, und die Klemmenspannung ist dieselbe. Der Arbeitspunkt liegt also auf \emph{beiden} Kennlinien. Zeichnen Sie die Quellengerade von $(0;\ \SI{12}{V})$ nach $(\SI{3}{A};\ 0)$ ein. Der Schnittpunkt mit der Kennlinie von N liegt bei
\[
  \ergebnis{I \approx \SI{1.5}{A}, \quad U \approx \SI{6}{V}}
\]
Kontrolle mit der Tabelle: Bei $U = \SI{6}{V}$ fließen durch N $\SI{1.50}{A}$; die Quelle liefert bei $\SI{1.5}{A}$ die Spannung $\SI{12}{V} - \SI{6}{V} = \SI{6}{V}$ \checkmark
\schritt{c) Widerstand im Arbeitspunkt}
\[
  R_\text{AP} = \frac{\SI{6}{V}}{\SI{1.5}{A}} = \ergebnis{\SI{4}{\ohm}}
\]
Kein Widerspruch: Dieser Wert gilt \emph{nur} in diesem einen Arbeitspunkt. Ändert sich die Quelle, ändert sich auch $R_\text{AP}$ -- man konnte ihn deshalb nicht vorab zur Berechnung nutzen.
\schritt{d) Höhere Leerlaufspannung}
Die neue Quellengerade verläuft \emph{parallel} zur alten (gleiche Steigung $-R_i$), von $(0;\ \SI{16}{V})$ nach $(\SI{4}{A};\ 0)$. Ablesen:
\[
  \ergebnis{I' \approx \SI{2.2}{A}, \quad U' \approx \SI{7.2}{V}}
\]
(Exakte Rechnung zur Kontrolle: Die Messwerte gehorchen $I = U^2/(\SI{24}{V^2/A})$. Gleichsetzen mit $I = (U_0'-U)/R_i$ führt auf $U^2 + 6U - 96 = 0$, also $U' = \SI{7.25}{V}$ und $I' = \SI{2.19}{A}$.)\\
Relative Stromänderung: $\dfrac{I'}{I} = \dfrac{2{,}19}{1{,}5} \approx 1{,}46$, also $\approx\ergebnis{+\SI{46}{\percent}}$, obwohl $U_0$ nur um $\SI{33}{\percent}$ gestiegen ist. Bei einer nichtlinearen Last sind Ursache und Wirkung nicht proportional!
\typfehler{Für d) den Widerstand $R_\text{AP} = \SI{4}{\ohm}$ aus c) weiterverwenden und $I' = \SI{16}{V}/\SI{8}{\ohm} = \SI{2}{A}$ rechnen -- das ist falsch, weil sich der Widerstand von N mit dem Arbeitspunkt ändert.}
\end{loesung}

\end{document}
